\documentclass[a4paper,11pt]{article}
\usepackage[czech]{babel}
\usepackage[T1]{fontenc}
\usepackage[utf8]{inputenc}
\usepackage{lmodern}
\usepackage{amsmath,amssymb}
\usepackage{geometry}
\usepackage{enumitem}
\usepackage{tikz}
\geometry{a4paper,left=17mm,right=17mm,top=13mm,bottom=14mm}
\pagestyle{empty}
\setlength{\parindent}{0pt}
\setlength{\parskip}{0pt}
\setlist[enumerate]{leftmargin=7mm,itemsep=0.5mm,topsep=1mm,parsep=0pt}
\newcommand{\TestHeader}[1]{%
\begin{center}
{\Large\bfseries ČTVRTLETNÍ TEST -- EXPONENCIÁLNÍ FUNKCE A ROVNICE}\\[1mm]
{\small\bfseries Včetně soustav exponenciálních rovnic}\hfill\textbf{Varianta #1}
\end{center}
\vspace{1mm}
Jméno a příjmení: \rule{8.5cm}{0.4pt}\hfill Třída: \rule{2cm}{0.4pt}\\[3mm]
Datum: \rule{3cm}{0.4pt}\hfill Počet bodů: \rule{2cm}{0.4pt}\hfill Známka: \rule{2cm}{0.4pt}\\[2mm]
\hrule\vspace{3mm}
}
\newcommand{\GrafSit}{%
\begin{tikzpicture}[x=0.48cm,y=0.48cm,>=stealth]
\draw[step=1,black!22,line width=0.23pt] (-6,-6) grid (6,6);
\draw[->,line width=.7pt] (-6.28,0)--(6.35,0) node[right] {$x$};
\draw[->,line width=.7pt] (0,-6.28)--(0,6.34) node[above] {$y$};
\foreach \x in {-6,-4,-2,2,4,6} {\draw (\x,.09)--(\x,-.09) node[below,font=\scriptsize] {\x};}
\foreach \y in {-6,-4,-2,2,4,6} {\draw (.09,\y)--(-.09,\y) node[left,font=\scriptsize] {\y};}
\node[below left,font=\scriptsize] at (0,0) {$0$};
\end{tikzpicture}}
% Zdroj typu grafu, definic a postupů: učebnice/02_Funkce_a_rovnice/
% Exponencialni_funkce_a_rovnice/Exponenciální funkce.tex a Exponenciální rovnice.tex.
% Soustavy jsou převzaty z úlohy 11: varianta A = a), varianta B = g).
% Výpočtové rovnice a transformační grafy jsou nové, sestavené pro tento test.

\begin{document}
\TestHeader{A}
\textbf{Teoretická část}
\begin{enumerate}
\item \textbf{Základ exponenciální funkce.}
\begin{enumerate}[label=\alph*),leftmargin=6mm,itemsep=0pt,topsep=0.6mm]
\item Jaká funkce vznikne z předpisu $f(x)=a^x$ při základu $a=1$?
\item Proč nemůže být základ exponenciální funkce záporný?
\end{enumerate}
\vspace{1.2mm}
\item \textbf{Monotónnost a vlastnosti grafu.}
\begin{enumerate}[label=\alph*),leftmargin=6mm,itemsep=0pt,topsep=0.6mm]
\item Je funkce $f(x)=a^{-x}$ rostoucí, nebo klesající? Proč? Uvažujte $a>0$, $a\neq1$.
\item Určete vodorovnou asymptotu a obor hodnot $g(x)=\left(\frac13\right)^{x-2}+4$.
\end{enumerate}
\vspace{1.2mm}
\item \textbf{Rovnice a substituce.}
\begin{enumerate}[label=\alph*),leftmargin=6mm,itemsep=0pt,topsep=0.6mm]
\item Jak zapíšete množinu řešení rovnice $5^{2-x}=0$ v $\mathbb R$?
\item Při substituci $t=2^x$ vyšly kořeny $t=1$ a $t=-\frac12$. Který vyloučíte, proč a jaký je výsledný kořen $x$?
\end{enumerate}
\vspace{1.2mm}
\end{enumerate}
\vspace{1mm}
\textbf{Početní část}\begin{enumerate}[start=4]
\item Načrtněte graf funkce $g(x)=2^{x-1}-2$. Určete její \textbf{definiční obor}, \textbf{obor hodnot} a \textbf{rovnici vodorovné asymptoty}.
\begin{center}\vspace{-1mm}\GrafSit\end{center}
\item Vyřešte exponenciální rovnici. Zapište i postup řešení.
\[
2^{2x+1}-9\cdot 2^x+4=0
\]
\vspace{20mm}
\item Vyřešte soustavu exponenciálních rovnic. Zapište i postup řešení.
\[
\begin{cases} 2^x+5\cdot 3^y=53,\\[2mm] 7\cdot 2^x-3^y=47. \end{cases}
\]
\vspace{1mm}
\end{enumerate}


\newpage
\TestHeader{B}
\textbf{Teoretická část}
\begin{enumerate}
\item \textbf{Základ exponenciální funkce.}
\begin{enumerate}[label=\alph*),leftmargin=6mm,itemsep=0pt,topsep=0.6mm]
\item Jaká funkce vznikne z $f(x)=a^{x-2}-3$ při $a=1$?
\item Proč je $f(x)=a^x$ při $a>0$, $a\neq1$ prostá? Zdůvodněte vlastností jejího grafu.
\end{enumerate}
\vspace{1.2mm}
\item \textbf{Monotónnost a vlastnosti grafu.}
\begin{enumerate}[label=\alph*),leftmargin=6mm,itemsep=0pt,topsep=0.6mm]
\item Je funkce $f(x)=a^{-2x}$ rostoucí, nebo klesající? Závislost na základu $a>0$, $a\neq1$ odůvodněte.
\item Určete vodorovnou asymptotu a obor hodnot $g(x)=\left(\frac14\right)^{x-2}-5$.
\end{enumerate}
\vspace{1.2mm}
\item \textbf{Rovnice a substituce.}
\begin{enumerate}[label=\alph*),leftmargin=6mm,itemsep=0pt,topsep=0.6mm]
\item Určete množinu řešení rovnice $7^{x+1}=-\frac17$ v $\mathbb R$ a zdůvodněte.
\item Při substituci $t=4^x$ vyšly kořeny $t=\frac14$ a $t=-2$. Vylučte nepřípustný kořen a určete $x$.
\end{enumerate}
\vspace{1.2mm}
\end{enumerate}
\vspace{1mm}
\textbf{Početní část}\begin{enumerate}[start=4]
\item Načrtněte graf funkce $g(x)=\left(\frac13\right)^{x+1}-3$. Určete její \textbf{definiční obor}, \textbf{obor hodnot} a \textbf{rovnici vodorovné asymptoty}.
\begin{center}\vspace{-1mm}\GrafSit\end{center}
\item Vyřešte exponenciální rovnici. Zapište i postup řešení.
\[
3^{2x+1}-10\cdot 3^x+3=0
\]
\vspace{20mm}
\item Vyřešte soustavu exponenciálních rovnic. Zapište i postup řešení.
\[
\begin{cases} 2\cdot 3^x+3\cdot 4^y=24,\\[2mm] 3\cdot 3^x-10\cdot 4^y=7. \end{cases}
\]
\vspace{1mm}
\end{enumerate}


\end{document}
